\pdfoutput=1
\pdfmapfile{+symbols.map}
\pdfmapfile{+cm.map}
\pdfmapfile{+cmextra.map}
\documentclass[10pt]{article}
\usepackage[margin=1in]{geometry}
\usepackage{amsmath,amssymb,amsthm}
\usepackage[hidelinks]{hyperref}
\newtheorem{theorem}{Theorem}
\title{The Smallest Cospectral Nonisomorphic Simple Graphs}
\author{Galaxy-0}
\date{October 9, 2026}
\begin{document}
\maketitle
\begin{abstract}
The minimum number of vertices in an adjacency-cospectral, nonisomorphic pair of finite simple undirected graphs is five. An explicit pair is the star $K_{1,4}$ and the disjoint union $C_4\sqcup K_1$. Both have characteristic polynomial $X^5-4X^3=X^3(X-2)(X+2)$ and spectrum $\{-2,0,0,0,2\}$. A self-contained Lean 4 project verifies the polynomial computation, nonisomorphism, and the exhaustive exclusion of all smaller pairs. This answers conjecture 00000003714 under the standard adjacency-spectrum convention.
\end{abstract}
\section{Scope and exact statement}
A graph here is finite, simple and undirected. Its spectrum is the multiset of eigenvalues of its adjacency matrix, including multiplicities. Equivalently, two graphs are cospectral when their adjacency characteristic polynomials agree. Disconnected graphs are allowed: the source imposes no connectedness hypothesis. Neither a Laplacian-spectrum assertion nor a minimum within connected graphs is claimed.
\begin{theorem}
There exist two nonisomorphic cospectral graphs on five vertices. If two graphs have the same number $n<5$ of vertices and have equal adjacency characteristic polynomials, they are isomorphic. Consequently five is the minimum possible order.
\end{theorem}
Graphs of different orders cannot be cospectral, since their characteristic polynomials are monic of different degrees.
\section{The explicit witnesses}
Use the common vertex set $\{0,1,2,3,4\}$. Define
\[
 E(S)=\{01,02,03,04\},\qquad E(H)=\{01,12,23,03\}.
\]
Thus $S=K_{1,4}$ and $H=C_4\sqcup K_1$, with vertex 4 isolated in $H$. Their adjacency matrices are
\[
 A_S=\begin{pmatrix}0&1&1&1&1\\1&0&0&0&0\\1&0&0&0&0\\1&0&0&0&0\\1&0&0&0&0\end{pmatrix},
 \qquad
 A_H=\begin{pmatrix}0&1&0&1&0\\1&0&1&0&0\\0&1&0&1&0\\1&0&1&0&0\\0&0&0&0&0\end{pmatrix}.
\]
Direct determinant expansion gives
\[
 \det(XI-A_S)=\det(XI-A_H)=X^5-4X^3=X^3(X-2)(X+2).
\]
This is also transparent from the star's two nonzero eigenvalues $\pm\sqrt4$ and the four-cycle's eigenvalues $2,0,0,-2$, with one extra zero from the isolated vertex. The Lean proof uses the determinant expansion itself, not these spectral facts as assumptions.

The graphs are not isomorphic: every vertex of $S$ has a neighbor, whereas vertex 4 of $H$ has none. If a bijective adjacency-preserving map existed, surjectivity would send some star vertex to 4, and its neighbor would contradict isolation. This argument is proved for arbitrary candidate bijections in Lean.
\section{Exhaustive proof of minimality}
For $n$ labeled vertices there are $2^{\binom n2}$ simple graphs. Encode each by the bits of an integer $g$ satisfying
\[
 0\le g<2^{n(n-1)/2}.
\]
The unordered pair $\{i,j\}$ with $i<j$ is assigned bit position $j(j-1)/2+i$. These positions run exactly once through $0,\ldots,\binom n2-1$. The diagonal is zero and adjacency is symmetric. Hence every finite simple graph on $n$ vertices occurs after a labeling.

The characteristic polynomial is computed from its definition:
\[
 \det(XI-A_g)=\sum_{\pi\in S_n}(-1)^{\operatorname{inv}(\pi)}
       \prod_{i=0}^{n-1}\bigl(X\delta_{i,\pi(i)}-(A_g)_{i,\pi(i)}\bigr).
\]
Integer coefficient lists in increasing degree order implement polynomial addition and convolution multiplication. Permutations are generated by inserting the first entry into every position of every recursively generated permutation of the remaining entries. Insertion generates each permutation once: removing that distinguished entry recovers its unique predecessor and insertion position. The Lean file additionally checks distinctness, lengths, entry bounds and factorial counts for all sizes used, $0\le n\le5$.

For each $n=0,1,2,3,4$, the checker considers \emph{every ordered pair} of graph masks. If the coefficient lists agree, it searches for a permutation witnessing an isomorphism. Every returned map is checked for injectivity, surjectivity, and preservation of every adjacency entry. The kernel-certified results are the theorems \texttt{small0} through \texttt{small4}. They cover respectively $1,1,2,8,64$ labeled graphs, hence $1,1,4,64,4096$ ordered pairs. There is no probabilistic sampling or dependence on a graph database.

The search is only a way to supply witnesses. Soundness does not assume that an external isomorphism program is correct: Lean proves that each selected map really is an isomorphism. If the search missed a required witness, the exhaustive theorem would fail to typecheck.
\section{Formal coverage and reproducibility}
\texttt{Main.lean} imports only \texttt{Std}. The main theorem \texttt{TLMC3714.conjecture3714} combines the explicit cospectral nonisomorphic pair with the universal result at every $n<5$. Supporting results verify the exact edge sets and the factorization of the characteristic polynomial.

The project is pinned to Lean 4.31.0. Run \texttt{lake build} from this directory. The verification log records the actual build and the output of \texttt{\#print axioms}. There are no proof placeholders, user-supplied axioms, native-decision oracles, or kernel bypasses. Finite certificates are established with \texttt{decide}. The final theorem depends only on the standard foundational axioms \texttt{propext} and \texttt{Quot.sound}.

The auxiliary Python program independently recomputes characteristic polynomials and the small-graph comparison as a readable cross-check. It is not a trusted input to the Lean proof. This submission establishes a minimum and an explicit attaining pair; it does not claim a classification of every larger cospectral pair.
\section{Source}
The exact bilingual conjecture is archived as \texttt{source.md}. The public source is \href{https://github.com/The-Last-Math-Competition/The-Last-Math-Competition/blob/main/conjectures/00000003714.md}{TLMC conjecture 00000003714}. All formal computations and definitions needed for this proof are included in the submission.
\end{document}
